\section{Cinemática}
\label{sec:cinematica}

\subsection{Por qué cuaterniones}

La dinámica se integra con el cuaternión y no con los ángulos de Euler. La
razón es práctica: la cinemática de Euler tiene una singularidad en
$\theta=\pm90^\circ$ (\emph{gimbal lock}), donde la matriz que convierte
velocidad angular en derivadas de los ángulos deja de ser invertible. Un
dron real no vuela de canto, pero un controlador mal sintonizado o una
maniobra acrobática pueden llevarlo ahí, y un simulador que se rompe
justo cuando el controlador falla no sirve para estudiar fallas. El
cuaternión no tiene singularidades, cuesta cuatro números en vez de nueve
(matriz de rotación) y su cinemática es bilineal, lo que facilita la
integración. Los ángulos de Euler se calculan solo para mostrar y para el
lazo de actitud.

\subsection{Ecuaciones}

La cinemática de traslación es la definición de velocidad:
\begin{equation}
\dot{\vect p} = \vect v .
\label{eq:pdot}
\end{equation}
La de rotación, con $\vect\omega$ expresada en el cuerpo
(\cod{math3d.q\_derivative}):
\begin{equation}
\dot{\vect q} = \tfrac12\,\vect q\otimes\begin{bmatrix}0\\ \vect\omega\end{bmatrix}
=\frac12\begin{bmatrix}
-q_x p - q_y q - q_z r\\
\phantom{-}q_w p + q_y r - q_z q\\
\phantom{-}q_w q - q_x r + q_z p\\
\phantom{-}q_w r + q_x q - q_y p
\end{bmatrix},
\label{eq:qdot}
\end{equation}
donde $\otimes$ es el producto de Hamilton. La matriz de rotación asociada
(\cod{math3d.q\_to\_rot}) es
\begin{equation}
\R(\vect q)=\begin{bmatrix}
1-2(q_y^2+q_z^2) & 2(q_xq_y-q_wq_z) & 2(q_xq_z+q_wq_y)\\
2(q_xq_y+q_wq_z) & 1-2(q_x^2+q_z^2) & 2(q_yq_z-q_wq_x)\\
2(q_xq_z-q_wq_y) & 2(q_yq_z+q_wq_x) & 1-2(q_x^2+q_y^2)
\end{bmatrix}.
\label{eq:Rq}
\end{equation}

\paragraph{Conversión con Euler.} De Euler a cuaternión
(\cod{euler\_to\_q}), con $c_\bullet=\cos(\bullet/2)$ y
$s_\bullet=\sin(\bullet/2)$:
\begin{equation}
\vect q=\begin{bmatrix}
c_\phi c_\theta c_\psi + s_\phi s_\theta s_\psi\\
s_\phi c_\theta c_\psi - c_\phi s_\theta s_\psi\\
c_\phi s_\theta c_\psi + s_\phi c_\theta s_\psi\\
c_\phi c_\theta s_\psi - s_\phi s_\theta c_\psi
\end{bmatrix}.
\end{equation}
Y de vuelta (\cod{q\_to\_euler}):
\begin{equation}
\begin{aligned}
\phi&=\operatorname{atan2}\!\big(2(q_wq_x+q_yq_z),\,1-2(q_x^2+q_y^2)\big),\\
\theta&=\arcsin\!\big(\sat_{[-1,1]}\,2(q_wq_y-q_zq_x)\big),\\
\psi&=\operatorname{atan2}\!\big(2(q_wq_z+q_xq_y),\,1-2(q_y^2+q_z^2)\big).
\end{aligned}
\end{equation}
La saturación del argumento del arcoseno evita que el redondeo produzca un
\texttt{NaN} cerca de $\theta=\pm90^\circ$.

\paragraph{Tasas de Euler y velocidad angular.} El lazo de ángulo produce
derivadas deseadas de los ángulos de Euler; el lazo de tasa necesita
velocidades angulares del cuerpo. La relación para Z-Y-X
(\cod{control/attitude.euler\_rates\_to\_body}) es
\begin{equation}
\vect\omega = W(\phi,\theta)\begin{bmatrix}\dot\phi\\ \dot\theta\\ \dot\psi\end{bmatrix},
\qquad
W=\begin{bmatrix}
1 & 0 & -\sin\theta\\
0 & \cos\phi & \sin\phi\cos\theta\\
0 & -\sin\phi & \cos\phi\cos\theta
\end{bmatrix}.
\label{eq:W}
\end{equation}
$W$ deja de ser invertible en $\cos\theta=0$: es la misma singularidad de
Euler, que por eso queda confinada al lazo de ángulo.

\paragraph{Normalización.} La integración numérica de \eqref{eq:qdot} no
conserva exactamente $\|\vect q\|=1$. Después de cada paso de física se
renormaliza (\cod{dynamics.renormalize\_quat}). Es una proyección de orden
del error de truncamiento, que no cambia el orden del método.

\subsection{Verificación}

El bloque B de \cod{model\_verification.py} toma 500 actitudes aleatorias
(roll y pitch hasta $\pm80^\circ$ y $\pm74^\circ$, yaw en $\pm172^\circ$) y
mide el peor caso de cada propiedad:

\begin{table}[H]
\centering
\small
\caption{Verificación numérica de la cinemática (peor caso en 500 actitudes).}
\label{tab:verif_cin}
\begin{tabularx}{\textwidth}{Xr}
\toprule
Propiedad comprobada & Error máximo \\
\midrule
Ortonormalidad, $\max|\R\T\R-I|$ y $|\det\R-1|$ & \val{eOrth} \\
$\R(\vect q)$ de \eqref{eq:Rq} contra $\R_z\R_y\R_x$ de \eqref{eq:zyx} & \val{eZYX} \\
Ida y vuelta Euler $\to\vect q\to$ Euler [rad] & \val{eRound} \\
$W$ de \eqref{eq:W} contra derivadas de Euler obtenidas por diferencias finitas de \eqref{eq:qdot} [rad/s] & \val{eW} \\
Jacobiano analítico $\partial(\R\T\vect d)/\partial\vect q$ usado por Madgwick y el EKF, contra diferencias finitas & \val{eJac} \\
\bottomrule
\end{tabularx}
\end{table}

Los dos primeros errores están en el nivel del redondeo de doble precisión;
los dos últimos son del orden del error de las diferencias finitas usadas
como referencia ($h=10^{-6}$ y $10^{-7}$), no del código.
